\documentclass[12pt]{article}
\usepackage[utf8]{inputenc}
\usepackage[T1]{fontenc}
\usepackage[margin=1in]{geometry}
\usepackage{setspace}
\usepackage{xcolor}
\usepackage[colorlinks=true, urlcolor=blue!60!black, linkcolor=blue!60!black]{hyperref}
\setlength{\parindent}{2em}
\setlength{\parskip}{1.25em}
\renewcommand{\baselinestretch}{1.0}
\begin{document}
\begin{center}
{\LARGE \textit{Trump v. Barbara}, 609 U.S. \_\_\_ (2026) (No. 25-365)\par}
\vspace{0.8em}
{\large Case Brief\par}
\end{center}
\vspace{1em}

\section{Memory Jogger}
In a 6--3 judgment (with a 5--4 constitutional majority), the Supreme Court invalidated President Trump's Executive Order No. 14160 and held that the Fourteenth Amendment's Citizenship Clause guarantees birthright citizenship to children born in the United States to parents who are unlawfully or temporarily present (\href{https://www.law.cornell.edu/supremecourt/text/25-365}{Cornell LII}; \href{https://www.courtlistener.com/opinion/10883542/trump-v-barbara-revisions-70126/}{CourtListener}).

\section{Detailed Case Facts}
On January 20, 2025, President Donald J. Trump issued Executive Order No. 14160, titled \textit{Protecting the Meaning and Value of American Citizenship}, 90 Fed. Reg. 8449 (2025), scheduled to take effect for children born on or after February 20, 2025. Section 1 of the Fourteenth Amendment provides: ``All persons born or naturalized in the United States, and subject to the jurisdiction thereof, are citizens of the United States and of the State wherein they reside.'' The Immigration and Nationality Act (INA), 8 U.S.C. \S\ 1401(a), codifies the same text (\href{https://www.law.cornell.edu/supremecourt/text/25-365}{Cornell LII}).

Executive Order No. 14160 declared that the privilege of U.S. citizenship does not automatically extend to persons born in the United States in two situations: (1) when the child's mother was unlawfully present in the United States and the father was not a U.S. citizen or lawful permanent resident (LPR) at the time of birth, or (2) when the child's mother's presence in the United States was lawful but temporary (such as on a tourist, student, or work visa) and the father was not a U.S. citizen or LPR at the time of birth. The Order asserted that such children are not ``subject to the jurisdiction'' of the United States under either the Fourteenth Amendment or the INA (\href{https://www.law.cornell.edu/supremecourt/text/25-365}{Cornell LII}).

Multiple federal district courts initially blocked the Order through universal injunctions, which reached the Supreme Court in \textit{Trump v. CASA, Inc.} (2025). On June 27, 2025, the Court ruled 6--3 in \textit{CASA} that federal district courts generally lack authority to issue universal injunctions, without deciding the underlying constitutionality of Executive Order No. 14160 (\href{https://en.wikipedia.org/wiki/Trump_v._Barbara}{Wikipedia}; \href{https://houstonlawreview.org/article/171991-from-universal-injunctions-to-nationwide-classes-what-_barbara-v-trump_-reveals-about-_casa_-s-remedial-framework.pdf}{Houston Law Review Online}).

That same day, June 27, 2025, the American Civil Liberties Union (ACLU) filed a putative class action, \textit{Barbara v. Trump}, No. 1:25-cv-00244, in the U.S. District Court for the District of New Hampshire. Several parents---including lead plaintiff ``Barbara,'' a Honduran national suing pseudonymously---filed suit on their own behalf and on behalf of (and in the name of) their children, arguing that Executive Order No. 14160 violated the Fourteenth Amendment's Citizenship Clause, 8 U.S.C. \S\ 1401(a), and the Administrative Procedure Act (\href{https://www.law.cornell.edu/supremecourt/text/25-365}{Cornell LII}; \href{https://clearinghouse.net/case/46748/}{Civil Rights Litigation Clearinghouse}).

\section{Procedural History}
\begin{itemize}
\item \textbf{U.S. District Court for the District of New Hampshire (\textit{Barbara v. Trump}, 790 F. Supp. 3d 80 (D.N.H. July 10, 2025)):} Judge Joseph N. Laplante narrowed the proposed class to children (excluding parents) born on or after February 20, 2025 who would be denied citizenship under the Order, provisionally certified that nationwide Rule 23(b)(2) class, and issued a class-wide preliminary injunction barring enforcement of Executive Order No. 14160. Judge Laplante stayed the injunction for seven days pending appeal; it took effect on July 17, 2025 (\href{https://clearinghouse.net/case/46748/}{Civil Rights Litigation Clearinghouse}; \href{https://houstonlawreview.org/article/171991-from-universal-injunctions-to-nationwide-classes-what-_barbara-v-trump_-reveals-about-_casa_-s-remedial-framework.pdf}{Houston Law Review Online}).
\item \textbf{U.S. Court of Appeals for the First Circuit (No. 25-1861):} The Government filed a notice of appeal on September 5, 2025, and petitioned the Supreme Court for a writ of certiorari before judgment on September 26, 2025 (\href{https://www.supremecourt.gov/DocketPDF/25/25-365/378052/20250926163053178_TrumpvBarbaraCertPet.pdf}{Certiorari Petition}).
\item \textbf{United States Supreme Court (609 U.S. \_\_\_ (2026), No. 25-365):} Granted certiorari before judgment in December 2025 (607 U.S. 1079 (2025)) on whether the Executive Order complies on its face with the Citizenship Clause and 8 U.S.C. \S\ 1401(a). Following oral argument on April 1, 2026 (with Solicitor General D. John Sauer arguing for Petitioners and ACLU National Legal Director Cecillia D. Wang for Respondents), the Court affirmed the District Court's judgment on June 30, 2026 (\href{https://www.scotusblog.com/cases/trump-v-barbara/}{SCOTUSblog}; \href{https://www.supremecourt.gov/opinions/25pdf/25-365_4hdj.pdf}{Supreme Court Slip Opinion}).
\end{itemize}

\section{Judicial Votes}
The Court affirmed the District Court's injunction by a 6--3 judgment, with a 5--4 majority on the constitutional question (\href{https://www.law.cornell.edu/supremecourt/text/25-365}{Cornell LII}; \href{https://www.scotusblog.com/2026/06/breaking-down-the-birthright-citizenship-decision/}{SCOTUSblog}):
\begin{itemize}
\item \textbf{Majority (5):} Chief Justice Roberts delivered the opinion of the Court, joined in full by Justices Sotomayor, Kagan, Barrett, and Jackson.
\item \textbf{Concurring (joining the majority in full):} Justice Jackson filed a concurring opinion, in which Justice Sotomayor joined as to the introduction and Part I.
\item \textbf{Concurring in the Judgment and Dissenting in Part (1):} Justice Kavanaugh filed an opinion concurring in the judgment (on statutory grounds under 8 U.S.C. \S\ 1401(a)) and dissenting from the Court's constitutional holding.
\item \textbf{Dissenting (3):} Justice Thomas filed the principal dissenting opinion, in which Justice Gorsuch joined. Justice Alito and Justice Gorsuch each also filed separate dissenting opinions.
\end{itemize}

\section{Holding}
\textbf{Judgment:} Affirmed. The District Court's class-wide preliminary injunction blocking enforcement of Executive Order No. 14160 was affirmed (\href{https://www.law.cornell.edu/supremecourt/text/25-365}{Cornell LII}).

\textbf{Rule Announced:} Children born in the United States to parents who are unlawfully or temporarily present in the country are ``subject to the jurisdiction'' of the United States and are therefore citizens at birth under the Citizenship Clause of the Fourteenth Amendment. ``Subject to the jurisdiction'' carries its ordinary territorial-sovereignty meaning---referring to the power of the United States to govern those within its territory---incorporating the common-law rule of \textit{jus soli} with only its narrow, settled exceptions (such as children of foreign diplomats, children born to hostile occupying forces, and historically members of quasi-sovereign Indian tribes). Neither the President by executive order nor Congress by statute may deny citizenship to persons guaranteed citizenship by the Clause (\href{https://www.law.cornell.edu/supremecourt/text/25-365}{Cornell LII}; \href{https://constitutioncenter.org/blog/supreme-court-strikes-down-trumps-birthright-citizenship-executive-order-in-landmark-decision}{National Constitution Center}).

\textbf{Holding vs. Dicta:} The binding constitutional holding (supported by five Justices) is that the Fourteenth Amendment itself confers birthright citizenship on children born on U.S. soil to unlawfully or temporarily present parents, rendering Executive Order No. 14160 facially unconstitutional. Because the majority resolved the case directly on the Fourteenth Amendment, it did not rely on an independent statutory holding under 8 U.S.C. \S\ 1401(a) (which Justice Kavanaugh alone treated as dispositive). The majority's discussion of the Civil Rights Act of 1866's phrasing (``not subject to any foreign power'') as raising ``more questions than answers'' and its historical commentary on post-1884 State Department rulings serve as supporting reasoning rather than independent holdings; meanwhile, the Court reaffirmed that \textit{United States v. Wong Kim Ark}'s repeated mentions of the parents' San Francisco domicile were factual descriptions rather than a constitutional limitation (\href{https://www.law.cornell.edu/supremecourt/text/25-365}{Cornell LII}).

\section{Analysis of Opinions}

\subsection{Majority Opinion (Chief Justice Roberts, joined by Justices Sotomayor, Kagan, Barrett, and Jackson)}
Chief Justice Roberts grounded the Court's opinion in four pillars: English and early American common law, the ordinary text of the Citizenship Clause, the repudiation of \textit{Dred Scott v. Sandford} (1857), and \textit{United States v. Wong Kim Ark} (1898). Under English common law---articulated in \textit{Calvin's Case} (1608) and Blackstone's \textit{Commentaries}---anyone born within the sovereign's dominions owed natural allegiance in return for protection and became a natural-born subject, regardless of how ``momentary and uncertain'' the parents' stay or whether the parents were subject to expulsion (such as Romani people). That \textit{jus soli} rule crossed the Atlantic, prevailed in every State after independence, was applied to carve out quasi-sovereign Indian tribes (citing Chancellor Kent in \textit{Goodell v. Jackson}), and was reaffirmed in \textit{Lynch v. Clarke} (N.Y. Ch. 1844) for children of temporary sojourners (\href{https://www.law.cornell.edu/supremecourt/text/25-365}{Cornell LII}).

After \textit{Dred Scott} abandoned \textit{jus soli} for a racial rule of bloodline, Attorney General Edward Bates's 1862 opinion, the Civil Rights Act of 1866, and the Fourteenth Amendment restored and permanently constitutionalized the common-law rule. Textually, ``jurisdiction'' in 1868 meant the sovereign's power to govern within its territory, as settled by Chief Justice Marshall in \textit{The Schooner Exchange v. McFaddon} (1812): foreign private visitors are fully ``amenable to the jurisdiction of the country,'' whereas narrow exceptions exist only to avoid degrading foreign sovereigns (diplomats) or where territorial jurisdiction was ceded to quasi-sovereign Indian tribes (\textit{Elk v. Wilkins}, 1884). Reaffirming \textit{Wong Kim Ark}, Roberts rejected the Government's and Justice Thomas's ``domicile'' and ``primary allegiance'' theories as ``dramatically revisionist,'' noting that the word ``domicile'' appeared only once in the Citizenship Clause debates and that words central to the Executive Order---``mother,'' ``father,'' ``lawful,'' ``temporary''---are absent from the constitutional text (\href{https://www.law.cornell.edu/supremecourt/text/25-365}{Cornell LII}).

\subsection{Concurring Opinion (Justice Jackson, joined by Justice Sotomayor as to Introduction and Part I)}
Joining the Court's opinion in full, Justice Jackson wrote separately to rebut Justice Thomas's framing of the Citizenship Clause as a narrow measure confined to freed slaves and those sharing their ``no other homeland'' status. Drawing on the antebellum Colored Conventions movement, Frederick Douglass, Abraham Lincoln, and the Reconstruction debates, she showed that Senator Lyman Trumbull deliberately withdrew an initial draft limited to ``persons of African descent'' in favor of universalist text---and that Congress knowingly extended citizenship to children of Chinese and Romani immigrants over Senator Edgar Cowan's objections and President Andrew Johnson's veto. In Part II (writing alone), Justice Jackson argued that Justice Thomas's group-specific reading sat uneasily with his usual insistence on a ``colorblind'' Constitution, maintaining that the Reconstruction Amendments enacted a universalist, ``anticaste, antisubordination reset'' for the Nation (\href{https://www.law.cornell.edu/supremecourt/text/25-365}{Cornell LII}).

\subsection{Concurring in the Judgment and Dissenting in Part (Justice Kavanaugh)}
Justice Kavanaugh agreed that Executive Order No. 14160 is unlawful, but would have decided the case solely on statutory grounds. Because Congress enacted 8 U.S.C. \S\ 1401(a) in 1940 and 1952 using the exact language of the Fourteenth Amendment after \textit{Wong Kim Ark} had read that language to recognize only four narrow exceptions, Congress incorporated \textit{Wong Kim Ark}'s rule by statute---a reading confirmed by 85 years of consistent Executive Branch practice. Thus, the President could not add new exceptions unilaterally. On the constitutional issue, however, Justice Kavanaugh dissented: rejecting the view that \textit{Wong Kim Ark}'s four exceptions are a ``closed set'' ``trapped in amber'' (citing \textit{United States v. Rahimi}, 2024), he argued that modern mass illegal immigration and global temporary travel are unanticipated post-1868 developments relevantly similar to the existing exceptions because the parents are non-citizens. Consequently, he concluded that Congress could constitutionally amend \S\ 1401(a) to exclude such children, though it had not done so (\href{https://www.law.cornell.edu/supremecourt/text/25-365}{Cornell LII}).

\subsection{Principal Dissenting Opinion (Justice Thomas, joined by Justice Gorsuch)}
In his longest dissent on the Court, Justice Thomas argued that the Citizenship Clause constitutionalized the Civil Rights Act of 1866's rule that birthright citizenship requires both birth on U.S. soil and \textit{domicile} (a permanent legal home), which confers ``full and complete'' political jurisdiction and primary allegiance. Rejecting \textit{Calvin's Case} as a ``feudal'' relic of perpetual servitude abandoned by the American settler ethos, he marshaled 19th-century prize cases (\textit{The Venus}, \textit{The Pizarro}), domicile and taxation precedents, floor statements by Senators Trumbull, Howard, Fessenden, and Williams, the \textit{Slaughter-House Cases} (1873), \textit{Elk v. Wilkins} (1884), late-19th-century State Department passport denials (e.g., Hausding, Greisser, de Bourry), and treatises by Thomas Cooley, Francis Wharton, and Henry Campbell Black. Reading \textit{Wong Kim Ark} as limited to children of permanently domiciled parents, he argued that because Executive Order No. 14160 had valid applications to temporary visitors and unlawfully present aliens lacking domicile, the facial challenge had to fail (\href{https://www.law.cornell.edu/supremecourt/text/25-365}{Cornell LII}; \href{https://www.scotusblog.com/2026/06/breaking-down-the-birthright-citizenship-decision/}{SCOTUSblog}).

\subsection{Dissenting Opinion (Justice Alito)}
Justice Alito called the decision ``one of the most important'' in the Court's history and ``a serious mistake.'' Focusing on allegiance and the Civil Rights Act's phrase ``not subject to any foreign power,'' he argued that the Fourteenth Amendment guarantees citizenship only to children who owe allegiance solely to the United States (not those automatically made nationals of their parents' home countries, except where parents have done everything in their power to become Americans, as in \textit{Wong Kim Ark}). Warning that the majority's rule rewards ``birth tourists'' and creates a ``powerful incentive to enter or remain in this country illegally,'' he urged that Congress---not the Fourteenth Amendment---should decide whether to grant statutory citizenship to children of undocumented or temporary parents (\href{https://www.theguardian.com/us-news/2026/jun/30/birthright-citizenship-explained-supreme-court}{The Guardian}; \href{https://forumtogether.org/article/explainer-supreme-court-reaffirms-birthright-citizenship-in-trump-v-barbara/}{National Immigration Forum}).

\subsection{Dissenting Opinion (Justice Gorsuch)}
In addition to joining Justice Thomas's dissent, Justice Gorsuch wrote a brief separate dissent emphasizing the demanding standard for facial challenges. Because, in his view, the historical record demonstrated that Executive Order No. 14160 was at least constitutional as applied to children born to temporary foreign visitors, the plaintiffs' facial challenge had to fail---leaving for another day harder as-applied questions, such as whether long-term unlawfully present residents might establish domicile sufficient for citizenship (\href{https://en.wikipedia.org/wiki/Trump_v._Barbara}{Wikipedia}; \href{https://www.scotusblog.com/2026/06/breaking-down-the-birthright-citizenship-decision/}{SCOTUSblog}).

\section{Examples: Future Applications}
\begin{enumerate}
\item \textbf{Same-Side (Citizenship Guaranteed):} \textit{Congressional Statute Excluding Children of Undocumented Immigrants.} Following Justice Kavanaugh's partial concurrence, Congress passes a bill amending 8 U.S.C. \S\ 1401(a) to deny birthright citizenship to children born in the United States whose parents entered without inspection. Under \textit{Trump v. Barbara}, this statute is unconstitutional: a five-Justice majority held that the Citizenship Clause itself---not merely the INA---guarantees citizenship to such children, placing the rule beyond ordinary legislation and requiring a constitutional amendment to alter.
\item \textbf{Same-Side (Citizenship Guaranteed):} \textit{State Refusal to Issue Birth Certificates to Children of Visa Holders.} A state legislature directs its vital-records agency to stamp ``Non-Citizen'' on birth certificates of infants born in state hospitals to mothers on B-2 tourist visas or F-1 student visas, denying them state citizenship and public-school eligibility. This falls squarely on the same side as \textit{Barbara}: temporary visitors are fully amenable to U.S. territorial jurisdiction under \textit{The Schooner Exchange} and \textit{Wong Kim Ark}, making their U.S.-born children both U.S. citizens and citizens of the State wherein they reside.
\item \textbf{Opposite-Side (No Constitutional Birthright Citizenship):} \textit{Child Born in Washington, D.C. to Accredited Foreign Ambassadors.} Two accredited ambassadors of a foreign nation stationed at their embassy in Washington, D.C. give birth to a daughter at a local hospital. Under \textit{Barbara}, \textit{Wong Kim Ark}, and \textit{Calvin's Case}, the child is \textit{not} a constitutional birthright citizen: diplomatic immunity and the common-law fiction of extraterritoriality exempt foreign ministers and their families from U.S. territorial jurisdiction to preserve sovereign dignity.
\item \textbf{Opposite-Side (No Constitutional Birthright Citizenship):} \textit{Child Born in a U.S. Territory Occupied by a Hostile Foreign Army.} During an armed conflict, a foreign military force seizes and occupies an American island, and a child is born there to members of the occupying force while U.S. governance is displaced. This falls within the classic common-law exception reaffirmed in \textit{Barbara} and \textit{Calvin's Case}: because the sovereign lacks actual territorial control and cannot extend protection in enemy-occupied areas, the child is not born ``subject to the jurisdiction'' of the United States.
\item \textbf{Fence-Sitter (Close / Unresolved Case):} \textit{Child Born in an Unincorporated U.S. Territory (e.g., American Samoa).} A child is born in American Samoa, an unincorporated U.S. territory where Congress has conferred non-citizen ``U.S. national'' status rather than statutory citizenship. While the child is unquestionably ``subject to the jurisdiction'' of the United States under \textit{Barbara}'s territorial-power test, the dispute turns on the Clause's \textit{other} textual element---whether unincorporated territories are ``in the United States'' under the \textit{Insular Cases} (e.g., \textit{Downes v. Bidwell}). \textit{Barbara} did not overrule the \textit{Insular Cases} (though Justice Jackson's concurrence listed them among regrettable decisions), leaving territorial birthright citizenship unresolved.
\end{enumerate}

\section{Critique}
\textbf{Originalist and Textualist Critiques:} \textit{Trump v. Barbara} sparked an intense intramural debate among originalists. Many originalist scholars---including amici such as Randy Barnett, Michael Ramsey, Evan Bernick, and the Cato Institute---supported the majority's conclusion that the original public meaning of ``subject to the jurisdiction thereof'' tracked territorial sovereignty and common-law \textit{jus soli}, rejecting domicile-based theories as ahistorical (\href{https://www.cato.org/legal-briefs/trump-v-barbara}{Cato Institute}; \href{https://scholarship.law.bu.edu/cgi/viewcontent.cgi?article=5201&context=faculty_scholarship}{BU Law Scholarship}). Conversely, conservative and originalist critics drawing on Kurt Lash and Ilan Wurman argued that the majority read the ``jurisdiction'' requirement out of the Clause, substituted medieval feudal subjecthood for republican consent and exclusive allegiance, and dismissed post-ratification Executive Branch practice (\href{https://journals.law.harvard.edu/jlpp/wp-content/uploads/sites/90/2026/09/Wessan-BirthrightCitizenshipSymposium-Final.pdf}{Harvard JLPP (2026)}). Post-liberal critics even claimed that the 5--4 split among self-described originalists exposed the methodological indeterminacy of originalism itself (\href{https://thenewdigest.substack.com/p/trump-v-barbara-and-the-end-of-originalism}{The New Digest}).

\textbf{Progressive and Doctrinal Perspectives:} Progressive constitutional scholars welcomed the result as a vital defense of Reconstruction's egalitarian core, while noting its fragility. Amanda Frost and Rachel E. Rosenbloom observed in the \textit{Harvard Law Review} that although Chief Justice Roberts's opinion preserved \textit{Wong Kim Ark}'s bright-line rule, four Justices (Thomas, Alito, Gorsuch, and Kavanaugh) signaled openness to congressional restrictions on birthright citizenship---meaning the constitutional guarantee held by a single vote. They urged defenders of birthright citizenship to build on Justice Jackson's anticaste, antisubordination concurrence rather than relying solely on common-law pedigree (\href{https://www.law.virginia.edu/node/2211356}{Frost \& Rosenbloom, 140 Harv. L. Rev. 137 (2026)}). Procedurally, scholars also highlighted how \textit{Barbara} demonstrated the durability of Rule 23(b)(2) nationwide class actions as a remedial backstop after \textit{Trump v. CASA, Inc.} curtailed universal injunctions (\href{https://houstonlawreview.org/article/171991-from-universal-injunctions-to-nationwide-classes-what-_barbara-v-trump_-reveals-about-_casa_-s-remedial-framework.pdf}{Houston Law Review Online}).

\textbf{Underlying Competing Values:} The opinions reflect a deep clash of constitutional values: (1) \textit{universal civic equality and anti-caste inclusion}---ensuring that no permanent hereditary underclass of U.S.-born non-citizens can emerge on American soil; (2) \textit{clarity and administrative certainty}---avoiding complex, manipulable inquiries into parental domicile or immigration status at birth; (3) \textit{consensual self-governance and immigration sovereignty}---the dissenters' concern that a political community should be able to define its own membership through democratic legislation; and (4) \textit{separation of powers}---highlighted by Justice Kavanaugh's insistence that even if the Constitution permits new exceptions, the President cannot override a century-old congressional statute by executive decree.

\section{Key Quotations}
\begin{quote}
``The Citizenship Clause uses jurisdiction in its ordinary sense---referring to the power of the United States to govern those within its territory.''
\end{quote}
\noindent --- Chief Justice Roberts, majority opinion (\href{https://www.law.cornell.edu/supremecourt/text/25-365}{Cornell LII}).

\begin{quote}
``Words appearing frequently in the Executive Order---`mother,' `father,' `lawful,' `temporary'---are absent from the Clause. For a simple reason: they did not matter.''
\end{quote}
\noindent --- Chief Justice Roberts, majority opinion (\href{https://www.law.cornell.edu/supremecourt/text/25-365}{Cornell LII}).

\begin{quote}
``Citizenship, then and now, was the right to have rights---to freely participate in our political community. The Framers of the Fourteenth Amendment extended that promise to `every free-born person in this land.' Cong. Globe, 39th Cong., 1st Sess., at 600 (Sen. Trumbull). We keep that promise today.''
\end{quote}
\noindent --- Chief Justice Roberts, majority opinion (\href{https://www.law.cornell.edu/supremecourt/text/25-365}{Cornell LII}).

\begin{quote}
``The Reconstruction Amendments were an anticaste, antisubordination reset for the Nation, not a mere spot treatment for the dark stain of slavery.''
\end{quote}
\noindent --- Justice Jackson, concurring (joined by Justice Sotomayor) (\href{https://www.law.cornell.edu/supremecourt/text/25-365}{Cornell LII}).

\begin{quote}
``In my view, the Executive Order does not violate the Fourteenth Amendment. But the Order does contravene a federal statute, 8 U. S. C. \S 1401(a). Congress could---consistent with the Fourteenth Amendment---amend \S 1401(a) or otherwise enact new legislation establishing exceptions to birthright citizenship for children born to foreign citizens unlawfully or temporarily in the country. But Congress has not yet done so.''
\end{quote}
\noindent --- Justice Kavanaugh, concurring in the judgment and dissenting in part (\href{https://www.law.cornell.edu/supremecourt/text/25-365}{Cornell LII}).


\end{document}
